% Part: many-valued-logic
% Chapter: three-valued-logics
% Section: kleene

\documentclass[../../../include/open-logic-section]{subfiles}

\begin{document}

\olfileid[fr]{mvl}{thr}{skl}

\olsection{Les logiques de Kleene}

Stephen Kleene a introduit deux
logiques à trois valeurs motivées
par une interprétation où les valeurs
de vérité sont les résultats de
procédures de calcul. Une procédure
peut fournir $\True$ ou $\False$,
mais elle peut aussi ne pas terminer.
Dans ce dernier cas, la valeur de
vérité correspondante est indéfinie,
ce que représente~$\Undef$.

Pour calculer la négation d'une
proposition~$!A$, on calculerait
d'abord la valeur de~$!A$, puis on
renverrait la valeur opposée. Si le
calcul de~$!A$ ne termine pas,
l'ensemble de la procédure ne termine
pas non plus : la négation de~$\Undef$
est donc~$\Undef$.

Pour calculer une conjonction
$!A\land !B$, deux approches sont
possibles. On peut d'abord calculer
$!A$, puis $!B$; le résultat est
$\True$ si les deux calculs donnent
$\True$, et $\False$ sinon. Si
l'un des calculs ne termine pas,
l'ensemble de la procédure ne termine
pas non plus. Dans cette approche,
une conjonction est donc indéfinie
dès que l'un de ses membres l'est;
il en va de même pour la disjonction.

Si, en revanche, nous pouvons évaluer
$!A$ et $!B$ en parallèle, nous
pouvons faire mieux. Dès qu'une des
deux procédures termine avec le
résultat~$\False$, nous pouvons
arrêter le calcul d'une conjonction :
son résultat doit être faux, même si
l'autre membre est indéfini. De même,
pour une disjonction calculée en
parallèle, nous pouvons nous arrêter
dès que l'une des procédures renvoie
$\True$ : le résultat doit être vrai.
Dans cette approche, nous pouvons
donc connaître la valeur d'un énoncé
composé même si celle d'un de ses
composants est indéfinie. Nous pouvons
alors lire $\Undef$ comme «~inconnue~»
plutôt que comme «~indéfinie~».

L'évaluation qui attend les deux résultats correspond aux tables
de la logique faible de Kleene. Celle qui peut s'arrêter dès
qu'un résultat suffit correspond aux tables de la logique forte.
Dans les deux cas, le conditionnel $!A\lif !B$ se définit par
$\lnot !A\lor !B$.

\begin{defn}
La \emph{logique forte de Kleene}
$\LogKs$ est définie par la matrice
suivante :
\begin{enumerate}
  \item Le fragment de $\Lang L_0$ dont les connecteurs sont $\lnot$,
  $\land$, $\lor$ et $\lif$.
  \item L'ensemble des valeurs de vérité
  $V=\{\True,\Undef,\False\}$.
  \item $\True$ est la seule valeur
  désignée, c'est-à-dire
  $V^+=\{\True\}$.
  \item Les fonctions de vérité sont
  données par les tables suivantes :
  \begin{center}
    \begin{tabular}{c|c}
      $\tf{\lnot}$ & \\
      \hline
      $\True$ & $\False$ \\
      $\Undef$ & $\Undef$ \\
      $\False$ & $\True$
    \end{tabular}
    \quad
    \begin{tabular}{c|ccc}
      $\tf{\land}[\LogKs]$ & $\True$ & $\Undef$ & $\False$ \\
      \hline
      $\True$ & $\True$ & $\Undef$ & $\False$ \\
      $\Undef$ & $\Undef$ & $\Undef$ & $\False$ \\
      $\False$ & $\False$ & $\False$ & $\False$
    \end{tabular}
    \\[2ex]
    \begin{tabular}{c|ccc}
      $\tf{\lor}[\LogKs]$ & $\True$ & $\Undef$ & $\False$ \\
      \hline
      $\True$ & $\True$ & $\True$ & $\True$ \\
      $\Undef$ & $\True$ & $\Undef$ & $\Undef$ \\
      $\False$ & $\True$ & $\Undef$ & $\False$
    \end{tabular}
    \quad
    \begin{tabular}{c|ccc}
      $\tf{\lif}[\LogKs]$ & $\True$ & $\Undef$ & $\False$ \\
      \hline
      $\True$ & $\True$ & $\Undef$ & $\False$ \\
      $\Undef$ & $\True$ & $\Undef$ & $\Undef$ \\
      $\False$ & $\True$ & $\True$ & $\True$
    \end{tabular}
  \end{center}
\end{enumerate}
\end{defn}

\begin{defn}
La \emph{logique faible de Kleene}
$\LogKw$ est définie par la matrice
suivante :
\begin{enumerate}
  \item Le fragment de $\Lang L_0$ dont les connecteurs sont $\lnot$,
  $\land$, $\lor$ et $\lif$.
  \item L'ensemble des valeurs de vérité
  $V=\{\True,\Undef,\False\}$.
  \item $\True$ est la seule valeur
  désignée, c'est-à-dire
  $V^+=\{\True\}$.
  \item Les fonctions de vérité sont
  données par les tables suivantes :
  \begin{center}
    \begin{tabular}{c|c}
      $\tf{\lnot}$ & \\
      \hline
      $\True$ & $\False$ \\
      $\Undef$ & $\Undef$ \\
      $\False$ & $\True$
    \end{tabular}
    \quad
    \begin{tabular}{c|ccc}
      $\tf{\land}[\LogKw]$ & $\True$ & $\Undef$ & $\False$ \\
      \hline
      $\True$ & $\True$ & $\Undef$ & $\False$ \\
      $\Undef$ & $\Undef$ & $\Undef$ & $\Undef$ \\
      $\False$ & $\False$ & $\Undef$ & $\False$
    \end{tabular}
    \\[2ex]
    \begin{tabular}{c|ccc}
      $\tf{\lor}[\LogKw]$ & $\True$ & $\Undef$ & $\False$ \\
      \hline
      $\True$ & $\True$ & $\Undef$ & $\True$ \\
      $\Undef$ & $\Undef$ & $\Undef$ & $\Undef$ \\
      $\False$ & $\True$ & $\Undef$ & $\False$
    \end{tabular}
    \quad
    \begin{tabular}{c|ccc}
      $\tf{\lif}[\LogKw]$ & $\True$ & $\Undef$ & $\False$ \\
      \hline
      $\True$ & $\True$ & $\Undef$ & $\False$ \\
      $\Undef$ & $\Undef$ & $\Undef$ & $\Undef$ \\
      $\False$ & $\True$ & $\Undef$ & $\True$
    \end{tabular}
  \end{center}
\end{enumerate}
\end{defn}

\begin{prop}
  $\LogKs$ et $\LogKw$ n'ont aucune tautologie dans le langage sans
  constantes des deux matrices.
\end{prop}

\begin{proof}
  Si $\pAssign{v}(p)=\Undef$ pour
  toute variable propositionnelle~$p$,
  alors toute formule~$!A$ a la
  valeur $\pValue{v}(!A)=\Undef$,
  car, dans les deux logiques,
  \[
    \tf{\lnot}(\Undef)
    =\tf{\lor}(\Undef,\Undef)
    =\tf{\land}(\Undef,\Undef)
    =\tf{\lif}(\Undef,\Undef)
    =\Undef.
  \]
  Or $\Undef\notin V^+$ pour
  $\LogKs$ comme pour $\LogKw$;
  $!A$ n'est donc pas désignée
  dans cette valuation.
\end{proof}

Bien que les logiques forte et faible
de Kleene n'aient aucune tautologie,
leurs relations de conséquence ne
sont pas triviales.

\begin{prob}
  Parmi les relations suivantes,
  lesquelles valent (a) dans la logique
  forte et (b) dans la logique faible
  de Kleene? Donner une table de vérité
  pour chacune.
  \begin{enumerate}
    \item $p,p\lif q\Entails q$;
    \item $p\lor q,\lnot p\Entails q$;
    \item $p\land q\Entails p$;
    \item $p\Entails p\land p$;
    \item $p\Entails p\lor q$.
  \end{enumerate}
\end{prob}

Dmitry Bochvar interprétait $\Undef$
comme «~dépourvu de sens~» et tenta
de s'en servir pour résoudre des
paradoxes tels que celui du Menteur,
en stipulant que les énoncés
paradoxaux prennent la valeur~$\Undef$.
Il introduisit une logique qui est
essentiellement la logique faible de
Kleene étendue par des connecteurs
supplémentaires, dont deux sont la
«~négation externe~» et l'opérateur
«~est indéfini~» :
\begin{center}
  \begin{tabular}{c|c}
    $\tf{\sim}$ & \\
    \hline
    $\True$ & $\False$ \\
    $\Undef$ & $\True$ \\
    $\False$ & $\True$
  \end{tabular}
  \quad
  \begin{tabular}{c|c}
    $\tf{+}$ & \\
    \hline
    $\True$ & $\False$ \\
    $\Undef$ & $\True$ \\
    $\False$ & $\False$
  \end{tabular}
\end{center}

\begin{prob}
  Peut-on définir $\sim$ dans la
  logique de Bochvar à partir de
  $\lnot$ et de~$+$? Autrement dit,
  peut-on trouver une formule contenant
  seulement la variable propositionnelle
  $p$, mais pas~$\sim$, et prenant
  toujours la même valeur de vérité
  que~$\mathord{\sim}p$? Donner une
  table de vérité pour justifier
  votre réponse.
\end{prob}

\end{document}
